\ifdefined\gkver\else\def\gkver{student}\fi
\documentclass[\gkver]{gaokaozhenti}
\begin{document}

\chapter{2015年重庆理综}
%% paperType: 理综物理部分

\begin{enumerate}
\item
%% number: 1
%% paperName: 2015年重庆理综第 1 题 6 分
%% typeId: 1
%% type: 单选题
%% chapter: 磁场
%% point: 带电粒子在磁场中的运动
%% method: 无
%% score: 6
%% degree: 650
%% duplicateId: 0
%% body:
图中曲线 a、b、c、d 为气泡室中某放射物发生衰变放出的部分粒子的径迹，气泡室中磁感应强度方向垂直于纸面向里。以下判断可能正确的是\xzanswer{D}

\onepicture[h=3.2cm]{figs/01.png}

\fourchoices[answer=D]
{a、b 为 $\beta$ 粒子的径迹}
{a、b 为 $\gamma$ 粒子的径迹}
{c、d 为 $\alpha$ 粒子的径迹}
{c、d 为 $\beta$ 粒子的径迹}

%% answer: D
%% memo:
\memoanswer{
	本题原卷及材料中未提供详解，待补充。
}

\item
%% number: 2
%% paperName: 2015年重庆理综第 2 题 6 分
%% typeId: 1
%% type: 单选题
%% chapter: 曲线运动
%% point: 人造地球卫星
%% method: 无
%% score: 6
%% degree: 650
%% duplicateId: 0
%% body:
宇航员王亚平在“天宫 1 号”飞船内进行了我国首次太空授课，演示了一些完全失重状态下的物理现象。若飞船质量为 $m$，距地面高度为 $h$，地球质量为 $M$，半径为 $R$，引力常量为 $G$，则飞船所在处的重力加速度大小为\xzanswer{B}

\fourchoices[answer=B]
{0}
{$\frac{GM}{(R+h)^{2}}$}
{$\frac{GMm}{(R+h)^{2}}$}
{$\frac{GM}{h^{2}}$}

%% answer: B
%% memo:
\memoanswer{
	本题原卷及材料中未提供详解，待补充。
}

\item
%% number: 3
%% paperName: 2015年重庆理综第 3 题 6 分
%% typeId: 1
%% type: 单选题
%% chapter: 运动和力
%% point: 加速减速模型
%% method: 无
%% score: 6
%% degree: 450
%% duplicateId: 0
%% body:
高空作业须安全带。如果质量为 $m$ 的高空作业人员不慎跌落，从开始跌落到安全带对人刚产生作用力前人下落的距离为 $h$（可视为自由落体运动）。此后经历时间 $t$ 安全带达到最大伸长，若在此过程中该作用力始终竖直向上，则该段时间安全带对人体的平均作用力大小为\xzanswer{A}

\fourchoices[answer=A]
{$\frac{m\sqrt{2gh}}{t}+mg$}
{$\frac{m\sqrt{2gh}}{t}-mg$}
{$\frac{m\sqrt{gh}}{t}+mg$}
{$\frac{m\sqrt{gh}}{t}-mg$}

%% answer: A
%% memo:
\memoanswer{
	本题原卷及材料中未提供详解，待补充。
}

\item
%% number: 4
%% paperName: 2015年重庆理综第 4 题 6 分
%% typeId: 1
%% type: 单选题
%% chapter: 电磁感应
%% point: 法拉第电磁感应定律
%% method: 无
%% score: 6
%% degree: 650
%% duplicateId: 0
%% body:
如图为无线充电技术中使用的受电线圈示意图，线圈匝数为 $n$，面积为 $S$，若在 $t_{1}$ 到 $t_{2}$ 时间内，匀强磁场平行于线圈轴线向右穿过线圈，其磁感应强度大小由 $B_{1}$ 均匀增加到 $B_{2}$，则该段时间线圈两端 a 和 b 之间的电势差 $\varphi_{1}-\varphi_{2}$\xzanswer{C}

\onepicture[h=3.2cm]{figs/04.png}

\fourchoices[answer=C]
{恒为 $\frac{nS(B_{2}-B_{1})}{t_{2}-t_{1}}$}
{从 0 均匀变化到 $\frac{nS(B_{2}-B_{1})}{t_{2}-t_{1}}$}
{恒为 $-\frac{nS(B_{2}-B_{1})}{t_{2}-t_{1}}$}
{从 0 均匀变化到 $-\frac{nS(B_{2}-B_{1})}{t_{2}-t_{1}}$}

%% answer: C
%% memo:
\memoanswer{
	本题原卷及材料中未提供详解，待补充。
}

\item
%% number: 5
%% paperName: 2015年重庆理综第 5 题 6 分
%% typeId: 1
%% type: 单选题
%% chapter: 运动和力
%% point: 牛顿定律中的图象
%% method: 无
%% score: 6
%% degree: 650
%% duplicateId: 0
%% body:
若货物随升降机运动的 $v-t$ 图像如图所示（竖直向上为正），则货物受到升降机的支持力 $F$ 与时间 $t$ 关系的图像可能是\xzanswer{B}

\onepicture[h=2.6cm]{figs/05.png}

\fourchoices[answer=B, ispicture=true, h=2.3cm]
{figs/05a.png}
{figs/05b.png}
{figs/05c.png}
{figs/05d.png}

%% answer: B
%% memo:
\memoanswer{
	本题原卷及材料中未提供详解，待补充。
}

\item
%% number: 6
%% paperName: 2015年重庆理综第 6 题 13 分
%% typeId: 4
%% type: 实验题
%% chapter: 电路
%% point: 伏安法测电阻
%% method: 无
%% score: 13
%% degree: 250
%% duplicateId: 0
%% body:
（1）（6 分）同学们利用如图所示方法估测反应时间。

首先，甲同学捏住直尺上端，使直尺保持竖直状态，直尺零刻度线位于乙同学的两指之间。当乙同学看见甲放开直尺时，立即用手指捏直尺，若捏住位置的刻度读数为 $x$，则乙同学的反应时间为\tkanswer[2.4cm]{$\sqrt{\frac{2x}{g}}$}（重力加速度为 $g$）。

基于上述原理，某同学用直尺制作测量反应时间的工具，若测量范围为 $0\sim0.4\Us$，则所用直尺的长度至少为\tkanswer[1.4cm]{80}$\Ucm$（$g$ 取 $10\Umsq$）；若以相等时间间隔在直尺上对应的长度是\tkanswer[1.8cm]{不相等}的（选填“相等”或“不相等”）。

\onepicture[h=4.6cm, align=right]{figs/06.png}

（2）（13 分）同学们测量某电阻丝的电阻 $R_{x}$，所用电流表的内阻与 $R_{x}$ 相当，电压表可视为理想电压表。

\begin{steps}
	\item
	若使用图 1 所示电路图进行实验，要使得 $R_{x}$ 的测量值更接近真实值，电压表的 a 端应连接到电路的\tkanswer[1.0cm]{c}点（选填“b”或“c”）；

	\item
	测得电阻丝的 $U-I$ 图如图 2 所示，则 $R_{x}$ 为\tkanswer[1.6cm]{4.1}$\UO$（保留两位有效数字）；

	\item
	实验中，随电压进一步增加电阻丝逐渐进入炽热状态。某同学发现对炽热电阻丝吹气，其阻值会变化。他们对此现象进行探究，在控制电阻丝两端的电压为 $10\UV$ 的条件下，得到电阻丝的电阻 $R_{x}$ 随风速 $v$（用风速计测）的变化关系如图 3 所示。由图可知当风速增加时，$R_{x}$ 会\tkanswer[1.4cm]{减小}（选填“增大”或“减小”）。在风速增加过程中，为保持电阻丝两端电压为 $10\UV$，需要将滑动变阻器 $R_{W}$ 的滑片向\tkanswer[1.0cm]{M}端移动（选填“M”或“N”）；

	\item
	为了通过电压表的示数来显示风速，同学们设计了如图 4 所示的电路。其中 $R$ 为两只阻值相同的电阻，$R_{x}$ 为两根相同的电阻丝，一根置于气流中，另一根不受气流影响，V 为待接入的理想电压表。如果要求在测量中，风速从零开始增加，电压表的示数也从零开始增加，则电压表的“$+$”和“$-$”端应分别连接到电路中的\tkanswer[1.0cm]{b}点和\tkanswer[1.0cm]{d}点（在“a”、“b”、“c”、“d”中选填）。
\end{steps}

\fourpicture[showlabel=false, hA=3.6cm, hB=3.6cm, hC=3.6cm, hD=3.6cm, align=right]
{figs/06a.png}
{figs/06b.png}
{figs/06c.png}
{figs/06d.png}

%% answer:
\jdanswer{
（1）$\sqrt{\frac{2x}{g}}$，$80\Ucm$，不相等；（2）①c；②$4.1\UO$；③减小、M；④b、d。
}
%% memo:
\memoanswer{
	本题原卷及材料中未提供详解，待补充。
}

\item
%% number: 7
%% paperName: 2015年重庆理综第 7 题 15 分
%% typeId: 6
%% type: 计算题
%% chapter: 磁场
%% point: 安培力
%% method: 无
%% score: 15
%% degree: 650
%% duplicateId: 0
%% body:
音圈电机是一种应用于硬盘、光驱等系统的特殊电动机，如图是某音圈电机的原理示意图，它由一对正对的磁极和一个正方形刚性线圈构成，线圈边长为 $L$，匝数为 $n$，磁极正对区域内的磁感应强度方向垂直于线圈平面竖直向下，大小为 $B$，区域外的磁场忽略不计。线圈左边始终在磁场外，右边始终在磁场内，前后两边在磁场内的长度始终相等。某时刻线圈中电流从 P 流向 Q，大小为 $I$。

\begin{enumerate}
	\item
	求此时线圈所受安培力的大小和方向。

	\item
	若此时线圈水平向右运动的速度大小为 $v$，求安培力的功率。
\end{enumerate}

\onepicture[h=3.8cm, align=right]{figs/07.png}

%% answer:
\jdanswer{
\begin{enumerate}
	\item
	$nIBL$，方向水平向右。

	\item
	$nIBLv$。
\end{enumerate}
}
%% memo:
\memoanswer{
	本题原卷及材料中未提供详解，待补充。
}

\item
%% number: 8
%% paperName: 2015年重庆理综第 8 题 16 分
%% typeId: 6
%% type: 计算题
%% chapter: 曲线运动
%% point: 平抛运动
%% method: 无
%% score: 16
%% degree: 450
%% duplicateId: 0
%% body:
同学们参照伽利略时期演示平抛运动的方法制作了如图所示的实验装置。图中水平放置的底板上竖直地固定有 M 板和 N 板。M 板上部有一半径为 $R$ 的 $\frac{1}{4}$ 圆弧形的粗糙轨道，P 为最高点，Q 为最低点，Q 点处的切线水平，距底板高为 $H$。N 板上固定有三个圆环。将质量为 $m$ 的小球从 P 处静止释放，小球运动至 Q 飞出后无阻碍地通过各圆环中心，落到底板上距 Q 水平距离为 $L$ 处。不考虑空气阻力，重力加速度为 $g$。求：

\begin{enumerate}
	\item
	距 Q 水平距离为 $\frac{L}{2}$ 的圆环中心到底板的高度；

	\item
	小球运动到 Q 点时速度的大小以及对轨道压力的大小和方向；

	\item
	摩擦力对小球做的功。
\end{enumerate}

\onepicture[h=4.6cm, align=right]{figs/08.png}

%% answer:
\jdanswer{
\begin{enumerate}
	\item
	$\frac{3}{4}H$

	\item
	$L\sqrt{\frac{g}{2H}}$，$mg\left(1+\frac{L^{2}}{2HR}\right)$，方向竖直向下。

	\item
	$mg\left(\frac{L^{2}}{4H}-R\right)$
\end{enumerate}
}
%% memo:
\memoanswer{
	本题原卷及材料中未提供详解，待补充。
}

\item
%% number: 9
%% paperName: 2015年重庆理综第 9 题 18 分
%% typeId: 6
%% type: 计算题
%% chapter: 磁场
%% point: 带电粒子在复合场中的运动
%% method: 无
%% score: 18
%% degree: 250
%% duplicateId: 0
%% body:
如图为某种离子加速器的设计方案。两个半圆形金属盒内存在相同的垂直于纸面向外的匀强磁场。其中 MN 和 M'N' 是间距为 $h$ 的两平行极板，其上分别有正对的两个小孔 O 和 O'，$O'N'=ON=d$，P 为靶点，$O'P=kd$（$k$ 为大于 1 的整数）。极板间存在方向向上的匀强电场，两极板间电压为 $U$，质量为 $m$、带电量为 $q$ 的正离子从 O 点由静止开始加速，经 O' 进入磁场区域。当离子打到极板上 O'N' 区域（含 N' 点）或外壳上时将会被吸收。两虚线之间的区域无电场和磁场存在，离子可匀速穿过，忽略相对论效应和离子所受的重力。求：

\begin{enumerate}
	\item
	离子经过电场仅加速一次后能打到 P 点所需的磁感应强度大小；

	\item
	能使离子打到 P 点的磁感应强度的所有可能值；

	\item
	打到 P 点的能量最大的离子在磁场中运动的时间和在电场中运动的时间。
\end{enumerate}

\onepicture[h=4.8cm, align=right]{figs/09.png}

%% answer:
\jdanswer{
\begin{enumerate}
	\item
	$\frac{2\sqrt{2qUm}}{qkd}$

	\item
	$\frac{2\sqrt{2nqUm}}{qkd}$（$n=1$，2，3，$\ldots$，$k^{2}-1$）

	\item
	在磁场中运动的时间为 $\frac{(2k^{2}-3)\pi mkd}{2\sqrt{2qUm(k^{2}-1)}}$，在电场中运动的时间为 $h\sqrt{\frac{2(k^{2}-1)m}{qU}}$。
\end{enumerate}
}
%% memo:
\memoanswer{
	本题原卷及材料中未提供详解，待补充。
}

\item
%% number: 10
%% paperName: 2015年重庆理综第 10 题 6 分
%% typeId: 6
%% type: 计算题
%% chapter: 气体性质
%% point: 查理定律
%% method: 无
%% score: 6
%% degree: 650
%% duplicateId: 0
%% body:
【选修 3—3】（12 分）

（1）（6 分）某驾驶员发现中午时车胎内的气压高于清晨时气压，且车胎体积增大。若这段时间胎内气体质量不变且可视为理想气体，那么\xzanswer{D}

\fourchoices[answer=D]
{外界对胎内气体做功，气体内能减小}
{外界对气体做功，气体内能增大}
{胎内气体对外做功，内能减小}
{胎内气体对外做功，内能增大}

（2）（6 分）北方某地的冬天室外气温很低，吹出的肥皂泡会很快冻结。若刚吹出时肥皂泡内气体温度为 $T_{1}$，压强为 $p_{1}$，肥皂泡冻结后泡内气体温度降为 $T_{2}$。整个过程中泡内气体视为理想气体，不计体积和质量变化，大气压为 $p_{0}$。求冻结后肥皂膜内外气体的压强差。

%% answer:
\jdanswer{
$\frac{T_{2}}{T_{1}}p-p_{0}$
}
%% memo:
\memoanswer{
	本题原卷及材料中未提供详解，待补充。
}

\item
%% number: 11
%% paperName: 2015年重庆理综第 11 题 6 分
%% typeId: 6
%% type: 计算题
%% chapter: 机械波
%% point: 波的图象
%% method: 无
%% score: 6
%% degree: 650
%% duplicateId: 0
%% body:
【选修 3—4】（12 分）

（1）（6 分）虹和霓是太阳光在水珠内分别经过一次和两次反射后出射形成的，可用白光照射玻璃球来说明。两束平行白光照射到透明玻璃球后，在水平的白色桌面上会形成 MN 和 PQ 两条彩色光带，光路如图所示，M、N、P、Q 点的颜色分别为\xzanswer{A}

\onepicture[h=3.6cm]{figs/11a.png}

\fourchoices[answer=A]
{紫、红、红、紫}
{红、紫、红、紫}
{红、紫、紫、红}
{紫、红、紫、红}

（2）（6 分）如图为一列沿 $x$ 轴正方向传播的简谐机械横波某时刻的波形图，质点 $P$ 的振动周期为 $0.4\Us$。求该波的波速并判断 P 点此时刻的振动方向。

\onepicture[h=3.8cm, align=right]{figs/11b.png}

%% answer:
\jdanswer{
该波的波速为 $2.5\Ums$，P 点此时的振动方向沿 $y$ 轴向下。
}
%% memo:
\memoanswer{
	本题原卷及材料中未提供详解，待补充。
}

\end{enumerate}

\end{document}
